\section{Formalization milestones and computational evidence}
\label{sec:formalization}

\subsection{A route to actual iPRA derivations}

The supplied Lean files are deliberately explicit about their status.
\path{FiniteSpecifications.lean} gives concrete numerical bounds, rational
predicates, a finite Davenport--Schmidt target, and several elementary ring
identities.  \texttt{ObjectTargets.lean} gives the scoped first-order shape of
St\o rmer completeness and applies the repository's existing transport to an
already supplied derivation.  Neither file supplies the main missing
Diophantine derivations.  The environment had no Lean executable, so the files
have not been elaborated or kernel checked.  They target the inspected
repository toolchain, \texttt{leanprover/lean4:v4.32.0}.

The following milestones specify the work required to complete the
formalization.  They are ordered so that each has a concrete, independently
reviewable output.

\begin{enumerate}
\item \textbf{Arithmetic definitions and characteristic predicates.}
Choose the numerical codec.  Give actual \texttt{PR} syntax trees for signed
arithmetic, rational operations, finite lists, bounded search, powers, gcd,
integer square root, and polynomial evaluation.  Prove their defining
equations and arithmetic invariants in \texttt{Prf PRAAx}.  Bound the
Euclidean algorithm by its initial input and then present it as a clocked
primitive recursion.  Compile every bounded quantifier into a Boolean fold.

\item \textbf{Finite counting and finite selection.}
Formalize the occupancy-table proof of pigeonhole, with a returned pair of
indices.  Formalize bounded minima, records, and prime-divisor search.  Each
search returns both a result and the finite certificate needed by the next
lemma; the proof of totality uses the exact bound from the article.

\item \textbf{The Roth fragment.}
Formalize rational reduction, the rational-root divisibility theorem, the
finite no-root test, polynomial differences, and bisection.  Prove the finite
separation inequality directly at precision $k$.  The cutoff then follows
by exponent arithmetic.  The object statement is conditional only on the
finite input checks and a positive integer $a-(d-2)b$.  Do not declare a
general effective-Roth axiom to fill the missing higher-degree result.

\item \textbf{St\o rmer.}
Build the exact ordered quadratic-pair arithmetic.  Prove both pigeonhole
constructions for Pell, the quotient's integrality and positivity, and its
explicit bound.  Prove bounded generation by the fundamental solution.
Formalize the binomial identity and its three prime-index cases.  Finish
with the coefficient enumeration, smoothness filter, completeness, cardinality,
and height bounds.  No external finiteness theorem is required.

\item \textbf{The finite Davenport--Schmidt core.}
Start with a rational working point and a finite local-separation predicate
through height $2B^2$.  Formalize the pigeonhole estimate, exact rational
minimal records, integral-plane lemma, independence bound, determinant
identity, and finite growth contradiction.  Prove the existence of a triple
inside the explicit box.  The two-by-two minor coordinate reduction and
the six-term determinant expansion are finite polynomial calculations.

\item \textbf{Root certificates and the represented-real schema.}
For a linear or quadratic polynomial, define its finite root code, reduce a
rational root, and prove the height comparison and distance estimate.  Add
the single-sample precision lemma and the error transfer.  State this as a
schema for each fixed \texttt{PR} name and separation modulus, together with
their arithmetic correctness premises.  Add reflection and reciprocal
normalization.  The algebraic specialization needs the finite Sylvester
adjugate identity and its coefficient bound.

\item \textbf{Assemble and transport.}
Construct the final \texttt{Prf PRAAx} derivations.  Apply
\texttt{pra\_theorem\_iff} to transfer them to the concise presentation.
Check the resulting declarations and their dependencies in Lean.  A
metatheoretic theorem about \texttt{Rat}, \texttt{Int}, or \texttt{Real} is
only an intermediate validation milestone.
\end{enumerate}

The finiteness bounds in the positive proofs are intentionally generous.
Their role is to make every primitive recursion and every induction invariant
explicit.  Optimizing them is a separate project and is not needed for the
claimed proof-theoretic result.

\subsection{Exact computations supplied with the report}

The Python reference code uses only integers, fractions, and integer square
roots.  It has no third-party dependency.  The verification run performs the
following genuinely different checks.

First, the St\o rmer enumerator is compared with independent multiplicative
generation of smooth integers up to $100000$, for the empty prime set and
the sets $\{3,5\}$, $\{2\}$, $\{2,3\}$, and $\{2,3,5\}$.  The complete
enumeration theorem is proved above; the finite comparison checks that the
implementation produces the expected values.  Pell identities and the key
smooth-coordinate property are also checked for every nonsquare
$2\leq D\leq50$, at the first ten powers of the fundamental solution.

For $S=\{2,3,5\}$, the result is exactly
\[
 1,2,3,4,5,8,9,15,24,80
\]
as the smaller entries of the neighbor pairs.  The JSON evidence lists every
Pell coefficient and every retained or rejected candidate.  The implementation
does not allocate $B(D)$: it decides $y<\mathrm{base}^{\mathrm{exponent}}$
by a bounded number of divisions of $y$ by the base.  This is an exact
comparison, not a smaller heuristic search cutoff.

Second, the quadratic Roth example uses an independent exact decision
procedure in $\mathbb Q(\sqrt2)$.  For $\delta=1/2$, the proved cutoff is
$q<36$.  The independent check through $q=72$ finds the reduced exceptional
fractions
\[
 1,\quad 2,\quad \frac32,\quad\frac75.
\]
Each has a separate rational bisection certificate.  Because the theorem
excludes every $q\geq36$ and the numerator range is exhaustive, these are
all the exceptional reduced fractions for this example.  The larger
computational range is a regression check, not the reason completeness holds.

Third, bounded searches find the following real quadratic approximants to
$\sqrt[3]{2}$.  The displayed decimals are explanatory rounded values; the
JSON file stores the exact rational isolating intervals and exact rational
upper bounds.

\begin{center}
\begin{tabular}{@{}rrrr@{}}
\toprule
Minimal polynomial & Height & Error upper bound & $H^3$ times error\\
\midrule
$6X^2-2X-7$ & 7 & $3.480\cdot10^{-4}$ & $0.120$\\
$13X^2-14X-3$ & 14 & $1.430\cdot10^{-4}$ & $0.393$\\
$27X^2-11X-29$ & 29 & $1.222\cdot10^{-5}$ & $0.298$\\
$59X^2-41X-42$ & 59 & $9.379\cdot10^{-7}$ & $0.193$\\
\bottomrule
\end{tabular}
\end{center}

Each polynomial is primitive with positive nonsquare discriminant; its
reported height is therefore its actual algebraic height.  Opposite rational
endpoint signs certify its real root, and a separate cubic interval certifies
the target.  These examples illustrate the exponent-three phenomenon; they
do not execute the enormous general height bound.

For this target, the algebraic modulus construction is concrete.  With
$F(X)=X^3-2$, $d=3$, $A=2$, and $R=2$, it gives
\[
 \rho(H)=648\max(2,H)^4.
\]
The polynomial is irreducible by the elementary Eisenstein argument at $2$,
and bisection on $[1,2]$ supplies its real root name.  Thus the theorem applies
with explicit arithmetic data, despite the intentionally impractical general
bound.

\subsection{Reproduction and honest completion criteria}

Run \texttt{python3 code/verify.py} from the archive directory to regenerate
the exact computational evidence.  Run \texttt{./build.sh} to compile the
LaTeX article twice.  The build uses ordinary \texttt{pdflatex} and standard
LaTeX packages.  The archived PDF has been compiled and visually inspected.

The file \texttt{lean/FORMALIZATION\_STATUS.txt} lists the exact unexecuted
Lean build commands and the pending object-theory work.  No completion claim
for a machine-checked iPRA proof is made.  Likewise, no mathematical claim is
made that full effective Roth has been resolved.  A future completion of
that part would need either a specified weaker constructive conclusion with
its full arithmetic proof, or a new effective-height theorem.  Merely
rephrasing classical finiteness or importing it as a hypothesis would not
complete that task.
